% Part: first-order-logic
% Chapter: completeness
% Section: construction-of-model

\documentclass[../../../include/open-logic-section]{subfiles}

\begin{document}

\iftag{FOL}
      {\olfileid{fol}{com}{mod}}
      {\olfileid{pl}{com}{mod}}

\olsection{Constructie van een model}

\begin{explain}
\iftag{FOL}{Voorlopig houden we ons niet bezig met~$\eq$: we willen
  alleen aantonen dat een consistente verzameling~$\Gamma$ van
  !!{sentence}s die~$\eq$ niet bevatten vervulbaar is. Eerst breiden we
  $\Gamma$ uit tot een consistente, !!{complete} en verzadigde
  verzameling~$\Gamma^*$. In dit geval is de definitie van een
  model~$\Struct{M(\Gamma^*)}$ eenvoudig. We nemen de verzameling gesloten
  termen van~$\Lang{L'}$ als domein. Aan elke !!{constant} kennen we
  zichzelf toe en, meer algemeen, zorgen we er voor elke gesloten term~$t$
  voor dat $\Value{t}{M(\Gamma^*)} = t$. Aan de !!{predicate}s kennen we
  extensies toe zodat een atomaire !!{sentence} waar is in
  $\Struct{M(\Gamma^*)}$ dan en slechts dan als zij tot~$\Gamma^*$ behoort.
  Daardoor zijn alle atomaire !!{sentence}s in~$\Gamma^*$ waar in
  $\Struct{M(\Gamma^*)}$. De overige zinnen zijn waar mits de
  verzameling~$\Gamma^*$ waarvan we uitgaan consistent, volledig en
  verzadigd is.}{We kunnen nu !!a{valuation} definiëren die alle
  $!A \in \Gamma$ waar maakt. Daartoe passen we eerst het lemma van
  Lindenbaum toe: we verkrijgen een volledige consistente
  $\Gamma^* \supseteq \Gamma$. De !!{propositional variable}s
  in~$\Gamma^*$ laten we~$\pAssign v(\Gamma^*)$ bepalen.}
\end{explain}

\iftag{FOL}{%
\begin{defn}[Termmodel]
\ollabel{defn:termmodel}
Laat $\Gamma^*$ !!a{complete}, consistente en verzadigde verzameling
!!{sentence}s zijn in een taal~$\Lang L$. Het \emph{termmodel}~$\Struct
M(\Gamma^*)$ van $\Gamma^*$ is de !!{structure} die als volgt wordt
gedefinieerd:
\begin{enumerate}
\item Het !!{domain}~$\Domain{M(\Gamma^*)}$ is de verzameling van alle
  gesloten termen van~$\Lang L$.
\item De interpretatie van !!a{constant} $c$ is $c$ zelf:
  $\Assign{c}{M(\Gamma^*)} = c$.
\item Aan de !!{function}~$f$ wordt de functie toegekend die bij de
  gesloten termen $t_1$, \dots, $t_n$ als argumenten de gesloten term
  $f(t_1, \dots, t_n)$ als waarde heeft:
\[
\Assign{f}{M(\Gamma^*)}(t_1, \dots, t_n) = f(t_1,\dots, t_n)
\]
\item Als $R$ een $n$-plaatsig !!{predicate} is, dan
  \[
  \tuple{t_1, \dots,
  t_n} \in \Assign{R}{M(\Gamma^*)} \text{ dan en slechts dan als } \Atom{R}{t_1, \dots,
    t_n} \in \Gamma^*.
  \]
\end{enumerate}
\end{defn}
}{%
\begin{defn}
\ollabel{defn:termmodel}
Veronderstel dat $\Gamma^*$ een volledige consistente verzameling
!!{formula}s is. Dan stellen we
\[
\pAssign v(\Gamma^*)(p) = \begin{cases}
  \True & \text{als $p \in \Gamma^*$}\\
  \False & \text{als $p \notin \Gamma^*$}
  \end{cases}
\]
\end{defn}
}

\iftag{FOL}{%
We gaan nu na dat inderdaad $\Value{t}{M(\Gamma^*)} = t$ geldt.

\begin{lem} 
\ollabel{lem:val-in-termmodel} Laat $\Struct M(\Gamma^*)$ het termmodel
uit \olref{defn:termmodel} zijn. Dan geldt $\Value{t}{M(\Gamma^*)} = t$.
\end{lem}
\begin{proof} 
 Het bewijs verloopt door inductie naar $t$. Het basisgeval, waarin $t$
 !!a{constant} is, volgt rechtstreeks uit de definitie van het termmodel.
 Veronderstel voor de inductiestap dat $t_1, \ldots, t_n$ gesloten termen
 zijn waarvoor $\Value{t_i}{M(\Gamma^*)} = t_i$ en dat $f$ een $n$-aire
 !!{function} is. Dan
\begin{align*}
\Value{f(t_1,\ldots,t_n)}{M(\Gamma^*)} &= \Assign{f}{M(\Gamma^*)}(\Value{t_1}
 {M(\Gamma^*)}, 
\ldots, \Value{t_n}{M(\Gamma^*)}) \\
&= \Assign{f}{M(\Gamma^*)}(t_1, \dots, t_n) \\
&= f(t_1,\dots, t_n),
\end{align*} 
en dus geldt de bewering door inductie voor elke gesloten term~$t$.
\end{proof}
}

\iftag{FOL}{%
\begin{explain}
  \iftag{prvEx}{!!^a{structure}~$\Struct{M}$ kan een existentieel
    gekwantificeerde !!{sentence}~$\lexists[x][!A(x)]$ waar maken zonder
    dat zij een instantie~$!A(t)$ waar maakt.}{}
  \iftag{prvAll}{!!^a{structure}~$\Struct{M}$ kan alle instanties~$!A(t)$
    van een universeel gekwantificeerde
    !!{sentence}~$\lforall[x][!A(x)]$ waar maken zonder
    $\lforall[x][!A(x)]$ waar te maken.}{} Dit komt doordat in het algemeen
  niet elk !!{element} van~$\Domain{M}$ de waarde van een gesloten term is
  ($\Struct{M}$ hoeft niet overdekt te zijn). Daarom is de
  vervulbaarheidsrelatie via variabelenbedelingen gedefinieerd. Voor ons
  termmodel~$\Struct{M(\Gamma^*)}$ zou dat echter niet nodig zijn, want dit
  model is overdekt. Dat is de inhoud van het volgende resultaat.
\end{explain}

\begin{prop}
\ollabel{prop:quant-termmodel}
Laat $\Struct M(\Gamma^*)$ het termmodel uit \olref{defn:termmodel} zijn.
\begin{tagenumerate}{prvEx,prvAll}
\tagitem{prvEx}{$\Sat{M(\Gamma^*)}{\lexists[x][!A(x)]}$ dan en slechts dan
  als $\Sat{M(\Gamma^*)}{!A(t)}$ voor ten minste één gesloten term~$t$.}{}
\tagitem{prvAll}{$\Sat{M(\Gamma^*)}{\lforall[x][!A(x)]}$ dan en slechts dan
  als $\Sat{M(\Gamma^*)}{!A(t)}$ voor alle gesloten termen~$t$.}{}
\end{tagenumerate}
\end{prop}

\begin{proof}
\begin{tagenumerate}{prvEx,prvAll}
\tagitem{prvEx}{%
  \iftag{probEx}{Opgave.}{Volgens \olref[syn][ass]{prop:sat-quant} geldt
    $\Sat{M(\Gamma^*)}{\lexists[x][!A(x)]}$ dan en slechts dan als voor
    ten minste één variabelenbedeling~$s$,
    $\Sat{M(\Gamma^*)}{!A(x)}[s]$. Omdat $\Domain{M(\Gamma^*)}$ uit de
    gesloten termen van~$\Lang{L}$ bestaat, geldt dit dan en slechts dan
    als er ten minste één gesloten term~$t$ is waarvoor $s(x) = t$ en
    $\Sat{M(\Gamma^*)}{!A(x)}[s]$. Volgens
    \olref[fol][syn][ext]{prop:ext-formulas},
    geldt $\Sat{M(\Gamma^*)}{!A(x)}[s]$ dan en slechts dan als
    $\Sat{M(\Gamma^*)}{!A(t)}[s]$, waarbij $s(x) = t$. Volgens
    \olref[fol][syn][ass]{prop:sentence-sat-true} geldt
    $\Sat{M(\Gamma^*)}{!A(t)}[s]$ dan en slechts dan als
    $\Sat{M(\Gamma^*)}{!A(t)}$, omdat $!A(t)$ een zin is.}}{}
\tagitem{prvAll}{%
  \iftag{probAll}{Opgave.}{Volgens \olref[syn][ass]{prop:sat-quant} geldt
    $\Sat{M(\Gamma^*)}{\lforall[x][!A(x)]}$ dan en slechts dan als voor
    elke variabelenbedeling $s$, $\Sat{M(\Gamma^*)}{!A(x)}[s]$. Bedenk
    dat $\Domain{M(\Gamma^*)}$ uit de gesloten termen van~$\Lang{L}$
    bestaat. Voor elke gesloten term~$t$ is $s(x) = t$ dus zo'n
    variabelenbedeling, en bij elke variabelenbedeling is $s(x)$ een
    zekere gesloten term~$t$. Volgens
    \olref[fol][syn][ext]{prop:ext-formulas} geldt
    $\Sat{M(\Gamma^*)}{!A(x)}[s]$ dan en slechts dan als
    $\Sat{M(\Gamma^*)}{!A(t)}[s]$, waarbij $s(x) = t$. Volgens
    \olref[fol][syn][ass]{prop:sentence-sat-true} geldt
    $\Sat{M(\Gamma^*)}{!A(t)}[s]$ dan en slechts dan als
    $\Sat{M(\Gamma^*)}{!A(t)}$, omdat $!A(t)$ een zin is.}}{}
\end{tagenumerate}
\end{proof}
}{}

\tagprob[FOL]{probEx,probAll}
\begin{prob}
Voltooi het bewijs van \olref[fol][com][mod]{prop:quant-termmodel}.
\end{prob}
\tagendprob

\begin{lem}[Waarheidslemma]
\ollabel{lem:truth} \iftag{FOL}{Veronderstel dat $!A$ geen~$\eq$ bevat. Dan}{}
geldt $\iftag{FOL}{\Sat{M(\Gamma^*)}}{\pSat{v(\Gamma^*)}}{!A}$ dan en slechts
dan als $!A \in \Gamma^*$.
\end{lem}

\begin{proof}
We bewijzen beide richtingen gelijktijdig door inductie naar $!A$.
\begin{enumerate}
\tagitem{prvFalse}{\indcase{!A}{\lfalse}
  {$\iftag{FOL}{\Sat/{M(\Gamma^*)}{\lfalse}}{\pSat/{v(\Gamma^*)}{\lfalse}}$
    volgens de definitie van vervulling. Daarentegen is $\lfalse \notin
    \Gamma^*$, omdat $\Gamma^*$ consistent is.}}{}

\tagitem{prvTrue}{\indcase{!A}{\ltrue}
  {$\iftag{FOL}{\Sat{M(\Gamma^*)}}{\pSat{v(\Gamma^*)}}{\ltrue}$
    volgens de definitie van vervulling. Anderzijds is $\ltrue \in
    \Gamma^*$, omdat $\Gamma^*$ consistent en !!{complete} is en
    $\Gamma^* \Proves \ltrue$.}}{}

\tagitem{FOL}{\indcase{!A}{R(t_1, \dots, t_n)}
  {$\Sat{M(\Gamma^*)}{\Atom{R}{t_1, \dots, t_n}}$ geldt dan en slechts dan
    als $\tuple{t_1, \dots, t_n} \in \Assign{R}{M(\Gamma^*)}$ (volgens de
    definitie van vervulling), en dat geldt dan en slechts dan als
    $R(t_1, \dots, t_n) \in \Gamma^*$ (volgens de constructie van
    $\Struct M(\Gamma^*)$).}}{\indcase{!A}{p}{$\pSat{v(\Gamma^*)}{p}$ geldt
    dan en slechts dan als $\pAssign v(\Gamma^*)(p) = \True$ (volgens de
    definitie van vervulling), en dat geldt dan en slechts dan als
    $p \in \Gamma^*$ (volgens de constructie van
    $\pAssign v(\Gamma^*)$).}}

\tagitem{prvNot}{%
  \indcase{!A}{\lnot !B}
    {$\iftag{FOL}{\Sat{M(\Gamma^*)}}{\pSat{v(\Gamma^*)}}{\indfrm}$ geldt
    dan en slechts dan als
    $\iftag{FOL}{\Sat/{M(\Gamma^*)}{!B}}{\pSat/{v(\Gamma^*)}{!B}}$
    (volgens de definitie van vervulling). Volgens de inductiehypothese
    geldt $\iftag{FOL}{\Sat/{M(\Gamma^*)}{!B}}{\pSat/{v(\Gamma^*)}{!B}}$
    dan en slechts dan als $!B \notin \Gamma^*$. Omdat
    $\Gamma^*$ consistent en !!{complete} is, geldt
    $!B \notin \Gamma^*$ dan en slechts dan als
    $\lnot !B \in \Gamma^*$.}}{}

\tagitem{prvAnd}{%
\iftag{probAnd}{%
  \indcase!{!A}{!B \land !C}{}}{%
  \indcase{!A}{!B \land
    !C}{$\iftag{FOL}{\Sat{M(\Gamma^*)}}{\pSat{v(\Gamma^*)}}{\indfrm}$
    geldt dan en slechts dan als zowel
    $\iftag{FOL}{\Sat{M(\Gamma^*)}}{\pSat{v(\Gamma^*)}}{!B}$ and
    $\iftag{FOL}{\Sat{M(\Gamma^*)}}{\pSat{v(\Gamma^*)}}{!C}$ gelden
    (volgens de definitie van vervulling), en dat geldt dan en slechts dan
    als zowel $!B \in \Gamma^*$ als $!C \in \Gamma^*$ (volgens de
    inductiehypothese). Volgens
    \olref[ccs]{prop:ccs}\olref[ccs]{prop:ccs-and} geldt dit dan en slechts
    dan als $(!B \land !C) \in \Gamma^*$.}}}{}

\tagitem{prvOr}{%
\iftag{probOr}{%
  \indcase!{!A}{!B \lor !C}{}}{%
  \indcase{!A}{!B \lor !C}
    {$\iftag{FOL}{\Sat{M(\Gamma^*)}}{\pSat{v(\Gamma^*)}}{\indfrm}$
     geldt dan en slechts dan als
     $\iftag{FOL}{\Sat{M(\Gamma^*)}}{\pSat{v(\Gamma^*)}}{!B}$ of
     $\iftag{FOL}{\Sat{M(\Gamma^*)}}{\pSat{v(\Gamma^*)}}{!C}$ geldt
     (volgens de definitie van vervulling). Volgens de inductiehypothese
     geldt dit dan en slechts dan als $!B \in \Gamma^*$ of
     $!C \in \Gamma^*$. Volgens
     \olref[ccs]{prop:ccs}\olref[ccs]{prop:ccs-or} geldt dat dan en slechts
     dan als $(!B \lor !C) \in \Gamma^*$.}}}{}
  
\tagitem{prvIf}{%
\iftag{probIf}{%
  \indcase!{!A}{!B \lif !C}{}}{%
  \indcase{!A}{!B \lif !C}
    {$\iftag{FOL}{\Sat{M(\Gamma^*)}}{\pSat{v(\Gamma^*)}}{\indfrm}$
     geldt dan en slechts dan als
     $\iftag{FOL}{\Sat/{M(\Gamma^*)}}{\pSat/{v(\Gamma^*)}}{!B}$ of
     $\iftag{FOL}{\Sat{M (\Gamma^*)}}{\pSat{v(\Gamma^*)}}{!C}$ geldt
     (volgens de definitie van vervulling). Volgens de inductiehypothese
     geldt dit dan en slechts dan als $!B \notin \Gamma^*$ of
     $!C \in \Gamma^*$. Volgens
     \olref[ccs]{prop:ccs}\olref[ccs]{prop:ccs-if} geldt dat dan en slechts
     dan als $(!B \lif !C) \in \Gamma^*$.}}}{}

\iftag{FOL}{%
\tagitem{prvAll}{%
  \iftag{probAll}{%
    \indcase!{!A}{\lforall[x][!B(x)]}{}}{%
    \indcase{!A}{\lforall[x][!B(x)]}{$\Sat{M(\Gamma^*)}{\indfrm}$ geldt
      dan en slechts dan als $\Sat{M(\Gamma^*)}{!B(t)}$ voor alle
      termen~$t$ (\olref{prop:quant-termmodel}). Volgens de
      inductiehypothese geldt dit dan en slechts dan als
      $!B(t) \in \Gamma^*$ voor alle termen~$t$. Volgens
      \olref[hen]{prop:saturated-instances} geldt dit op zijn beurt dan en
      slechts dan als $\lforall[x][!A(x)] \in \Gamma^*$.}}}{}

\tagitem{prvEx}{%
  \iftag{probEx}{%
    \indcase!{!A}{\lexists[x][!B(x)]}{}}{%
    \indcase{!A}{\lexists[x][!B(x)]}{$\Sat{M(\Gamma^*)}{\indfrm}$ geldt
      dan en slechts dan als $\Sat{M(\Gamma^*)}{!B(t)}$ voor ten minste
      één term~$t$ (\olref{prop:quant-termmodel}). Volgens de
      inductiehypothese geldt dit dan en slechts dan als
      $!B(t) \in \Gamma^*$ voor ten minste één term~$t$. Volgens
      \olref[hen]{prop:saturated-instances} geldt dit op zijn beurt dan en
      slechts dan als $\lexists[x][!B(x)] \in \Gamma^*$.}}}{}
}{}
\end{enumerate}
\end{proof}


\tagprob[FOL]{probOr,probAnd,probIf,probEx,probAll}

\begin{prob}
Voltooi het bewijs van \olref[fol][com][mod]{lem:truth}.
\end{prob}

\tagendprob

\tagprob[notFOL]{probOr,probAnd,probIf}

\begin{prob}
Voltooi het bewijs van \olref[pl][com][mod]{lem:truth}.
\end{prob}

\tagendprob

\end{document}
